% === Begin Label Data ===
% ID: 120
% 难度星级: 
% 标签: 
% 备注: 
% 组卷引用次数: 0
% === End  Label Data ===

\begin{problem}{2021}{G}{新高考I卷}{16}{数列}
某校学生在研究民间剪纸艺术时，发现剪纸时经常会沿纸的某条对称轴把纸对折，规格为 $20dm \times 12dm$ 的长方形纸，对折 1 次共可以得到 $10dm \times 12dm, 20dm \times 6dm$ 两种规格的图形，它们的面积之和 $S_1=240dm^2$，对折 2 次共可以得到 $5dm \times 12dm, 10dm \times 6dm, 20dm \times 3dm$ 三种规格的图形，它们的面积之和 $S_2=180dm^2$，以此类推，则对折 4 次共可以得到不同规格图形的种数为\underline{\hspace{4em}}；如果对折 $n$ 次，那么 $\sum\limits_{k=1}^n S_k=$\underline{\hspace{4em}} $dm^2$.
\end{problem}